\chapter{Logika intuicjonistyczna}\label{ch:ncl-intuitionistic}

\section{Motywacja i interpretacja konstruktywna}

W logice klasycznej zdanie jest prawdziwe albo fałszywe
w danym wartościowaniu. Logika intuicjonistyczna
rozróżnia natomiast posiadanie dowodu zdania od
niemożności uzyskania jego dowodu. Nie utożsamiamy
braku dowodu $\varphi$ z dowodem $\neg\varphi$.
Nie jest to kolejna trójwartościowa tablica spójników
w sensie Definicji~\ref{def:ncl-matrix}: informację
o dowodach można rozszerzać, a znaczenie implikacji
zależy od sposobu takiego rozszerzania.

\begin{definicja}[Język zdaniowy]
\label{def:ncl-ipc-language}
Formuły języka intuicjonistycznego rachunku zdań
(IPC) budujemy z ustalonego nieskończonego zbioru
zmiennych zdaniowych, stałej $\bot$ oraz spójników
$\wedge,\vee,\to$. Przyjmujemy skróty
\[
 \neg\varphi:=\varphi\to\bot,\qquad
 \top:=\bot\to\bot,\qquad
 \varphi\leftrightarrow\psi:=
 (\varphi\to\psi)\wedge(\psi\to\varphi).
\]
\end{definicja}

\subsection{Interpretacja Brouwera--Heytinga--Kołmogorowa}
Interpretacja BHK opisuje konstrukcje uzasadniające
poszczególne formuły. Jest to motywacja dla reguł
dowodowych, nie przyjęta tutaj formalna definicja
relacji konsekwencji.
\begin{itemize}
\item Dowód $\varphi\wedge\psi$ składa się z dowodu
      $\varphi$ i dowodu $\psi$.
\item Dowód $\varphi\vee\psi$ wskazuje, który człon
      alternatywy został dowiedziony, oraz zawiera jego dowód.
\item Dowód $\varphi\to\psi$ jest konstrukcją
      przekształcającą dowolny dowód $\varphi$
      w dowód $\psi$.
\item $\bot$ nie ma konstruktywnego uzasadnienia.
      Dowód $\neg\varphi$ przekształca hipotetyczny
      dowód $\varphi$ w dowód sprzeczności.
\end{itemize}
W szczególności dowód $\neg\neg\varphi$ nie dostarcza
bezpośrednio dowodu $\varphi$: pozwala odrzucić
każdą hipotetyczną konstrukcję dowodu $\neg\varphi$.
Formalną różnicę między tymi zdaniami pokażemy
za pomocą semantyki Kripkego.

\begin{przyklad}[Konstruktywna alternatywa]
\label{ex:ncl-bhk-disjunction}
Niech $\varphi$ wyraża własność obiektu, której
nie rozstrzygnęliśmy. Sam fakt, że nie mamy dowodu
$\varphi$, nie daje konstrukcji dowodu $\neg\varphi$.
Nie dostarcza więc również dowodu
$\varphi\vee\neg\varphi$ w interpretacji BHK.
Jeśli jednak mamy procedurę rozstrzygającą daną
własność wraz z uzasadnieniem wyniku, jej zastosowanie
może dostarczyć dowodu tej konkretnej alternatywy.
Nie oznacza to, że logika intuicjonistyczna odrzuca
wszystkie instancje prawa wyłączonego środka.
\end{przyklad}

\section{Dedukcja naturalna}

\begin{definicja}[Intuicjonistyczna dedukcja naturalna]
\label{def:ncl-ipc-natural-deduction}
Przez $\Gamma\vdash_{\mathrm I}\varphi$ rozumiemy
istnienie skończonego wyprowadzenia $\varphi$,
którego otwarte założenia należą do zbioru $\Gamma$.
Reguła założenia pozwala użyć każdej formuły
z aktualnego kontekstu. Reguły wprowadzania
i eliminacji spójników mają postać
\[
 \frac{\varphi\qquad\psi}{\varphi\wedge\psi}
 \ (\wedge I),\qquad
 \frac{\varphi\wedge\psi}{\varphi}\ (\wedge E_1),
 \qquad
 \frac{\varphi\wedge\psi}{\psi}\ (\wedge E_2),
\]
\[
 \frac{\varphi}{\varphi\vee\psi}\ (\vee I_1),
 \qquad
 \frac{\psi}{\varphi\vee\psi}\ (\vee I_2),
\]
\[
 \frac{\varphi\vee\psi\qquad
       \begin{matrix}[\varphi]^i\\\vdots\\\theta\end{matrix}
       \qquad
       \begin{matrix}[\psi]^j\\\vdots\\\theta\end{matrix}}
      {\theta}\ (\vee E_{i,j}),
\]
\[
 \frac{\begin{matrix}[\varphi]^i\\\vdots\\\psi\end{matrix}}
      {\varphi\to\psi}\ (\to I_i),
 \qquad
 \frac{\varphi\to\psi\qquad\varphi}{\psi}\ (\to E),
 \qquad
 \frac{\bot}{\varphi}\ (\bot E).
\]
Nawiasy kwadratowe oznaczają założenia tymczasowe,
a indeksy identyfikują wystąpienia zamykane przez
daną regułę. Wolno zamknąć zero wystąpień założenia.
W regule $\vee E$ zamykamy tylko wskazane założenia
w odpowiednich dwóch podwyprowadzeniach.
Pozostałe otwarte założenia pozostają otwarte.
\end{definicja}
Reguła $\bot E$ zachowuje zasadę eksplozji.
Usunięcie tej reguły prowadzi do logiki minimalnej,
której nie utożsamiamy z IPC.
Nie przyjmujemy klasycznej reguły pozwalającej
wnioskować $\varphi$ z $\neg\neg\varphi$.

\begin{stwierdzenie}[Własności strukturalne dedukcji]
\label{prop:ncl-ipc-structural}
Relacja $\vdash_{\mathrm I}$ jest zwrotna względem
założeń, monotoniczna i niezmiennicza względem
jednolitych podstawień. Spełnia również regułę:
\[
 \Gamma\vdash_{\mathrm I}\delta\quad(\delta\in\Delta),
 \qquad \Delta\vdash_{\mathrm I}\varphi
 \quad\Longrightarrow\quad
 \Gamma\vdash_{\mathrm I}\varphi.
\]
Każde wyprowadzenie wykorzystuje skończenie wiele
otwartych założeń.
\end{stwierdzenie}
\begin{proof}
Zwrotność wynika z reguły założenia, a monotoniczność
z tego, że powiększenie zbioru dozwolonych otwartych
założeń nie zmienia wyprowadzenia.
Przy jednolitym podstawieniu zastępujemy wszystkie
formuły ich obrazami, zachowując strukturę drzewa
oraz indeksy zamykanych wystąpień założeń.
Każdy schemat reguły pozostaje instancją tej samej reguły.
W szczególności $\bot$ pozostaje stałą.

Dla ostatniej reguły weźmy skończone wyprowadzenie
$\varphi$ z $\Delta$. Każdy jego otwarty liść
oznaczony $\delta$ zastępujemy kopią wyprowadzenia
$\delta$ z $\Gamma$. Przed podstawieniem zmieniamy
indeksy zamykanych założeń na świeże, aby żadna
reguła wyprowadzenia zewnętrznego nie zamknęła
przypadkowo otwartego założenia dołączonej kopii.
Otrzymujemy skończone drzewo, którego otwarte
założenia należą do $\Gamma$.
Skończoność wykorzystywanych założeń wynika
ze skończonej liczby liści drzewa.
\end{proof}

\begin{twierdzenie}[Twierdzenie o dedukcji]
\label{thm:ncl-ipc-deduction}
Dla dowolnego zbioru formuł $\Gamma$ zachodzi
\[
 \Gamma\cup\{\varphi\}\vdash_{\mathrm I}\psi
 \iff \Gamma\vdash_{\mathrm I}\varphi\to\psi.
\]
\end{twierdzenie}
\begin{proof}
Jeżeli istnieje wyprowadzenie $\psi$ z
$\Gamma\cup\{\varphi\}$, oznaczamy wszystkie jego
otwarte wystąpienia założenia $\varphi$ świeżym
indeksem i zamykamy je regułą $\to I$.
Pozostałe otwarte założenia należą do $\Gamma$.
Jeżeli nie ma wystąpień $\varphi$, stosujemy
wprowadzenie implikacji z pustym zbiorem zamykanych
wystąpień.

Odwrotnie do wyprowadzenia $\varphi\to\psi$ z
$\Gamma$ dołączamy założenie $\varphi$ i stosujemy
$\to E$. Otrzymane wyprowadzenie ma otwarte
założenia w $\Gamma\cup\{\varphi\}$.
\end{proof}

\begin{przyklad}[Podstawowe wyprowadzenia]
\label{ex:ncl-ipc-basic-proofs}
W IPC zachodzą między innymi
\[
 \vdash_{\mathrm I}\varphi\to\varphi,\qquad
 \vdash_{\mathrm I}\varphi\to(\psi\to\varphi),
 \qquad
 \vdash_{\mathrm I}\varphi\to\neg\neg\varphi.
\]
Pierwsze wyprowadzenie otrzymujemy, zamykając
założenie $\varphi$. W drugim zaczynamy od
$\varphi$, wprowadzamy $\psi\to\varphi$ bez
używania założenia $\psi$, a następnie zamykamy
$\varphi$. Trzecie ma postać
\[
 \frac{
   \frac{
     \frac{[\varphi]^1\qquad[\varphi\to\bot]^2}
          {\bot}\ (\to E)
   }{(\varphi\to\bot)\to\bot}\ (\to I_2)
 }{\varphi\to((\varphi\to\bot)\to\bot)}\ (\to I_1).
\]
Zewnętrzna implikacja zamyka założenie $\varphi$,
a wewnętrzna założenie $\neg\varphi$.
\end{przyklad}

\begin{stwierdzenie}[Potrójna negacja i osłabiona alternatywa]
\label{prop:ncl-ipc-triple-negation}
W IPC zachodzi
\[
 \vdash_{\mathrm I}\neg\neg\neg\varphi
       \leftrightarrow\neg\varphi,\qquad
 \vdash_{\mathrm I}\neg\neg(\varphi\vee\neg\varphi).
\]
\end{stwierdzenie}
\begin{proof}
Kierunek $\neg\varphi\to\neg\neg\neg\varphi$
jest instancją poprzedniego przykładu.
Dla drugiego kierunku zakładamy
$\neg\neg\neg\varphi$ i $\varphi$.
Z $\varphi$ uzyskujemy $\neg\neg\varphi$,
a przez eliminację implikacji otrzymujemy $\bot$.
Zamykając kolejno założenia $\varphi$ oraz
$\neg\neg\neg\varphi$, dostajemy żądaną implikację.
Łączymy obie implikacje regułą $\wedge I$.

Niech teraz $h=\neg(\varphi\vee\neg\varphi)$
będzie założeniem. Z dodatkowego założenia
$\varphi$ otrzymujemy $\varphi\vee\neg\varphi$
przez $\vee I_1$, a z $h$ otrzymujemy $\bot$.
Zamykając dodatkowe założenie, uzyskujemy
$\neg\varphi$ przy nadal otwartym $h$.
Stosując $\vee I_2$ oraz $h$, otrzymujemy ponownie
$\bot$. Zamykamy $h$, co daje
$\neg\neg(\varphi\vee\neg\varphi)$.
\end{proof}

\begin{stwierdzenie}[Intuicjonistyczne prawa de Morgana]
\label{prop:ncl-ipc-demorgan}
W IPC zachodzą
\[
 \vdash_{\mathrm I}\neg(\varphi\vee\psi)
       \leftrightarrow(\neg\varphi\wedge\neg\psi),
 \qquad
 \vdash_{\mathrm I}(\neg\varphi\vee\neg\psi)
       \to\neg(\varphi\wedge\psi).
\]
\end{stwierdzenie}
\begin{proof}
Z założenia $\neg(\varphi\vee\psi)$ oraz
dodatkowego założenia $\varphi$ uzyskujemy
$\varphi\vee\psi$ i $\bot$. Zamknięcie $\varphi$
daje $\neg\varphi$. Analogicznie otrzymujemy
$\neg\psi$. Wprowadzamy koniunkcję i zamykamy
pierwotne założenie.
Odwrotnie z $\neg\varphi\wedge\neg\psi$
uzyskujemy oba człony. Z dodatkowego założenia
$\varphi\vee\psi$ przez $\vee E$ uzyskujemy
$\bot$: w pierwszym podwyprowadzeniu stosujemy
$\neg\varphi$ do $\varphi$, w drugim
$\neg\psi$ do $\psi$. Zamykamy kolejno
założenie alternatywy i założenie koniunkcji.

Dla drugiej tezy zakładamy $\neg\varphi\vee\neg\psi$
oraz $\varphi\wedge\psi$. Z koniunkcji
uzyskujemy $\varphi$ i $\psi$.
Eliminacja alternatywy daje $\bot$ w obu przypadkach:
$\neg\varphi$ w połączeniu z $\varphi$ oraz
$\neg\psi$ w połączeniu z $\psi$.
Zamykamy oba pierwotne założenia w odpowiedniej
kolejności, otrzymując żądaną implikację.
\end{proof}
Implikacji odwrotnej w drugim prawie de Morgana
nie można ogólnie wyprowadzić w IPC. Kontrmodel
podamy po wprowadzeniu semantyki Kripkego.
